% ch9_conclusion.tex — 第九章 结论
\section{结论}

\subsection{研究发现总结}

本研究采用大语言模型（DeepSeek-V3.2）从亚马逊平台1{,}331条评论中提取基于Fishbein多属性态度理论的消费者态度指数（$\text{ATT} = \sum b_i \times e_i$），并通过四个假设系统验证其预测效度。

H1/H2分析（$N=1{,}157$）的核心发现：ATT与Rating显著正相关（$r=0.722$，$p<0.001$），有序Logistic回归证实ATT显著正向预测Rating（$\beta=1.867$，$z=19.8$，$p<0.001$），OLS $R^2=0.560$，H1获得支持，表明LLM提取的多属性态度具有良好的预测效度。评论者经验对ATT--Rating关系存在显著负向调节（$\beta=-0.008$，$p=0.001$），低经验评论者的ATT--Rating斜率（0.741）高于高经验评论者（0.573），H2以探索性结果获得支持。这一负向调节方向揭示了高经验评论者在评分时可能引入竞品比较等多维度权衡，使态度与评分的简单对应关系趋于复杂，挑战了"经验线性增强态度--行为一致性"的传统假设。

H3/H4分析（$N=298$）的核心发现：ATT显著正向预测BI（$\beta=0.167$，$p=0.004$），H3获得支持；在控制Rating后ATT仍提供显著增量预测力（$\Delta R^2=0.018$，$F=8.32$，$p=0.004$），H4获得支持。需要注意，H3/H4分析存在同源方差（Harman单因子$=76.8\%$）和选择偏差（$t=4.604$，$p<0.001$）两项局限，结论仅适用于"明确表达行为意向的评论者"子群，不宜推广至全体评论者。

\subsection{理论贡献与管理启示}

\textbf{方法论贡献。}本研究提出并验证了基于LLM的多属性态度提取方法（ATT计算准确率100\%，重测信度$r=0.892$，$e_i$一致率97.2\%），为"Text as Data"研究范式\cite{grimmer2022}在消费者行为研究中的应用提供了方法论基础。三数据源设计（LLM文本提取／平台客观评分／全品类行为数据）从根源上消除了CMV问题\cite{podsakoff2003}，为LLM文本分析研究中的CMV控制提供了可复用范式。H3/H4分析中ATT与BI同源于同一文本（Harman单因子$=76.8\%$）的CMV局限与H1/H2分析的三数据源设计形成鲜明对比，进一步凸显了独立数据源设计的方法论价值。本研究进一步探索了以重复购买行为或有用票数作为BI外部代理的可能性，从数据结构（重复评论比例在不同规模产品中均稳定低于1\%）和构念效度（假阴性率极高、重复评论$\neq$重复购买、意向--行为鸿沟）两个层面系统排除了上述替代方案，从而反证了三数据源设计的不可替代性。

\textbf{理论贡献。}H4证明ATT在控制Rating之后仍对BI具有显著增量预测力（$\Delta R^2=0.018$），回应了"ATT只是Rating代理变量"的质疑，表明多属性分解能够捕捉单一评分所压缩的信念结构信息，支持了Fishbein多属性态度理论\cite{fishbein1975}的区分效度。H2的负向调节发现丰富了消费者专业性理论\cite{mitchell2022}在在线评论情境下的内涵，提示研究者在利用LLM提取的态度变量预测评分行为时，需要将评论者经验纳入调节机制的考量。

\textbf{管理启示。}本研究验证了从评论文本中自动提取多属性态度的可行性，为企业提供了低成本的消费者态度监测工具。H4表明多属性分解提供了单一评分之外的额外信息，企业不应仅依赖评分数据评估产品表现，而应深入分析消费者对各属性的信念结构，识别"高重要性+低满意度"的改进维度。H2提示针对高经验评论者群体，其评分行为涉及更复杂的比较性权衡，企业在解读该群体的评分时应结合多属性信念数据进行综合分析。

\subsection{研究局限与未来展望}

本研究存在以下主要局限。第一，单一产品外部效度局限：本研究仅以伯特小蜜蜂护手霜套装为研究对象，方法的普适性需在更广泛的产品类别和电商平台中验证。第二，H3/H4分析的CMV与选择偏差：ATT与BI同源于同一文本（Harman单因子$=76.8\%$），且BI存在组ATT显著更高（$t=4.604$，$p<0.001$），两项局限共同制约了H3/H4结论的推广范围。第三，$b_i$评估的跨模型稳定性：信念强度ICC$=0.249$低于可接受阈值，不同LLM在0--3量表的数值校准上缺乏统一参照，影响了方法的可复现性。第四，reviewer\_experience的代理效度：以品类评论总数衡量消费者专业性存在局限，其与评论质量指标（$r$均接近0）的弱相关提示该代理指标有待改进。第五，H2调节方向的解释边界：H2的负向调节方向虽以探索性假设获得支持，但本研究对其机制的解释（比较性评价框架、评分策略保守化）属于事后推演，无法区分"负向调节确实反映了经验导致评价复杂化"与"反向结果部分源于代理指标局限或样本结构"两种可能。未来研究需通过更直接的消费者专业性测量手段和预先注册的假设对这一调节方向进行独立验证。

未来研究可沿三个方向推进：其一，将LLM多属性态度提取方法扩展至多产品类别、多语言和多平台，检验方法的稳健性；其二，采用独立数据源或实验设计改进H3/H4分析的CMV控制，并通过倾向得分匹配等方法缓解选择偏差；其三，探索基于规则或混合模型的$b_i$赋值方案，提高信念强度评估的跨模型一致性，同时在Fishbein理论框架内整合主观规范、感知行为控制等变量，构建更完整的行为预测模型。
